% v11 (R02) RQ5: a stronger locator inside the safety workflow (WP I, protocol I). Numbers: numbers-wpi macros.
\section{RQ5: A Stronger Locator Inside the Safety Workflow}\label{sec:locator}
Sections~\ref{sec:external} and~\ref{sec:student} show that the cascade is not the best locator. Its change labels, two of its three guards and its issue states, however, do not depend on how a window was found. The indentation guard and the copy guard inspect only the returned window; the variable-height window belongs to the cascade's own search and has no counterpart over another locator. This section tests that separation directly: can the safety-oriented workflow improve the reliability of an already strong locator? We placed CodeMapper \cite{hu2026codemapper}, the most accurate locator in our comparisons, and the contextual tracker inside the workflow, and then a language model. Each locator returns a window or no match. The window passes through the quotation comparison that assigns the cascade's labels (\emph{labels only}: \lbl{content-match} or \lbl{moved} when the normalized quotation is found unchanged, \lbl{changed} otherwise, \lbl{removed} on no match) and then through the indentation guard and the corrected copy guard, IG+AGi. The guards never move a window, so placement stays the locator's own and only the requests change. Because CodeMapper sees one file pair at a time, we also evaluated it with the cascade's cross-file search as a fallback when it reports no match.

The protocol, adapters and hypotheses were recorded with code hashes before any guarded output was computed (protocol~I in the artifact), and each amendment was recorded before the outcomes it concerns. The guarded locators were evaluated on the real histories, the student-authored repositories, the human-annotated benchmark and a new stress seed generated after the protocol was fixed, which holds out \wiFrN{} indentation, duplicate, short-anchor and compound cases. AGi was specified after the student and compound results of an earlier copy-guard revision had been read, so for AGi only the new seed is held out. Table~\ref{tab:locator} gives the results.

% Generated by tools/v11/build_wpi_tex.py from summary-20261102T183407.json; do not edit.
\begin{table*}[t]
\caption{Locators With and Without the Safety Layer (\%)}\label{tab:locator}
\centering\scriptsize
\setlength{\tabcolsep}{1.5pt}
\begin{tabular}{@{}lrrrrrr@{\hspace{3pt}}rrrrrr@{\hspace{3pt}}rrrrrr@{\hspace{3pt}}rrrrrr@{}}\toprule
 & \multicolumn{6}{c}{Real histories} & \multicolumn{6}{c}{Student repositories} & \multicolumn{6}{c}{Human-annotated} & \multicolumn{6}{c}{Held-out stress} \\
 & \multicolumn{6}{c}{$n{=}$15,064} & \multicolumn{6}{c}{$n{=}$7,882} & \multicolumn{6}{c}{$n{=}$200} & \multicolumn{6}{c}{$n{=}$23,550} \\
\cmidrule(lr){2-7}\cmidrule(lr){8-13}\cmidrule(lr){14-19}\cmidrule(lr){20-25}
Method & Place & Miss & Silent & Misdir & Lost & Extra & Place & Miss & Silent & Misdir & Lost & Extra & Place & Miss & Silent & Misdir & Lost & Extra & Place & Miss & Silent & Misdir & Lost & Extra \\ \midrule
\multicolumn{25}{@{}l}{\emph{Cascade}} \\
\quad deployed & 91.8 & 1.3 & 0.3 & 1.6 & 6.4 & 0.0 & 91.8 & 3.1 & 1.5 & 1.7 & 5.5 & 0.0 & 62.2 & 8.5 & 3.5 & 9.0 & 27.6 & 34.3 & 66.0 & 71.8 & 37.1 & 0.9 & 3.3 & 0.9 \\
\quad +IG+AGi+VW & 93.7 & 0.1 & 0.0 & 1.0 & 5.4 & 0.7 & 93.2 & 0.2 & 0.4 & 2.0 & 4.9 & 3.8 & 67.3 & 3.8 & 0.5 & 9.0 & 26.3 & 44.3 & 69.7 & 4.9 & 5.2 & 29.9 & 3.0 & 32.8 \\
\quad +IG+AGi+VW, $\theta{=}0.55$ & 97.8 & 0.0 & 0.0 & 2.1 & 0.4 & 0.7 & 96.6 & 0.1 & 0.4 & 3.0 & 0.7 & 3.8 & 87.2 & 3.8 & 0.5 & 21.0 & 1.9 & 44.3 & 69.6 & 4.7 & 5.0 & 32.6 & 1.4 & 34.0 \\
\multicolumn{25}{@{}l}{\emph{CodeMapper}} \\
\quad labels only & 97.9 & 0.4 & 0.1 & 1.5 & 0.5 & 0.8 & 92.4 & 0.5 & 0.2 & 3.6 & 3.6 & 6.1 & 93.6 & 3.1 & 0.0 & 7.5 & 1.3 & 35.7 & 83.0 & 19.3 & 7.0 & 4.6 & 10.6 & 13.4 \\
\quad +IG+AG & 97.9 & 0.0 & 0.0 & 1.6 & 0.5 & 1.9 & 92.4 & 0.1 & 0.1 & 3.6 & 3.6 & 8.3 & 93.6 & 0.8 & 0.0 & 7.5 & 1.3 & 50.0 & 83.0 & 2.0 & 3.0 & 8.6 & 10.6 & 38.0 \\
\quad +IG+AGi & 97.9 & 0.0 & 0.0 & 1.6 & 0.5 & 1.3 & 92.4 & 0.1 & 0.1 & 3.6 & 3.6 & 7.4 & 93.6 & 2.3 & 0.0 & 7.5 & 1.3 & 44.3 & 83.0 & 3.1 & 3.7 & 7.9 & 10.6 & 37.0 \\
\multicolumn{25}{@{}l}{\emph{CodeMapper, cross-file fallback}} \\
\quad labels only & 98.3 & 0.4 & 0.1 & 1.5 & 0.1 & 0.3 & 95.2 & 1.1 & 0.8 & 3.6 & 0.3 & 2.0 & 93.6 & 3.1 & 0.0 & 7.5 & 1.3 & 35.7 & 85.7 & 22.3 & 8.7 & 4.6 & 7.2 & 8.8 \\
\quad +IG+AGi$^{*}$ & 98.3 & 0.0 & 0.0 & 1.6 & 0.1 & 1.0 & 95.2 & 0.1 & 0.1 & 4.2 & 0.3 & 4.9 & 93.6 & 2.3 & 0.0 & 7.5 & 1.3 & 44.3 & 85.7 & 4.7 & 4.6 & 8.7 & 7.2 & 33.3 \\
\multicolumn{25}{@{}l}{\emph{Contextual tracker}} \\
\quad labels only & 95.6 & 0.3 & 0.1 & 0.6 & 2.9 & 1.7 & 91.9 & 0.5 & 0.2 & 1.2 & 7.2 & 7.2 & 82.7 & 3.8 & 0.0 & 3.0 & 14.7 & 35.7 & 82.0 & 22.0 & 7.7 & 0.1 & 11.8 & 13.2 \\
\quad +IG+AGi & 95.6 & 0.0 & 0.0 & 0.7 & 2.9 & 2.1 & 91.9 & 0.1 & 0.1 & 1.2 & 7.2 & 8.9 & 82.7 & 3.1 & 0.0 & 3.0 & 14.7 & 44.3 & 82.0 & 3.3 & 4.1 & 3.7 & 11.8 & 37.9 \\
\bottomrule
\end{tabular}
\tablenote{}{Place: correct start among targets that survive. Miss: changed or deleted targets not returned for re-check. Silent: targets placed on wrong code without a request, over all targets. Misdir: targets placed on wrong code with a request, over all targets, including deleted code relocated onto other code. Lost: surviving targets with no location. Extra: requests on targets whose content is preserved. \emph{Labels only} passes a locator's window through the quotation comparison that assigns the label; IG and AGi never move a window, so the other locators' placement is unchanged, whereas the cascade's guarded rows also include VW, which changes window height and placement. Every locator sees the anchor's original path only (no Git rename hint). Real histories and student repositories are Python; the human-annotated regions span several languages and are seen as file pairs, so cross-file recovery is not exercised there. Held-out stress is a fresh seed of the indentation, duplicate, short-anchor and compound suites generated after the protocol was fixed. The pinned CodeMapper is not fully deterministic (Section~\ref{sec:prefixes}); its guarded rows reuse the windows of its unguarded rows, so the two are paired. $^{*}$Composed after outcomes were read (exploratory).}
\end{table*}


\textbf{The safety layer made the strong locator safer.} With labels only, CodeMapper left \wiFrCmQMiss\% of the needed re-checks unrequested on the held-out seed and placed \wiFrCmQSilent\% of the targets on wrong code without a request. With IG+AGi these fell to \wiFrCmGiMiss\% and \wiFrCmGiSilent\% (paired cluster interval of the change in missed re-checks \wiFrCmGiVsCmQMissLo{} to \wiFrCmGiVsCmQMissHi{} points). On the student repositories missed re-checks fell from \wiStuCmQMissK{} to \wiStuCmGiMissK{} of \wiStuCmQMissN{} and silent misplacements from \wiStuCmQSilentK{} to \wiStuCmGiSilentK{}; on the real histories both fell from \wiRealCmQMissK{} to \wiRealCmGiMissK{}; on the human-annotated regions missed re-checks fell from \wiHumCmQMissK{} to \wiHumCmGiMissK{} of \wiHumCmQMissN{}. The contextual tracker responded in the same way. The gain has two sources. IG catches indentation changes that leave the normalized text intact, which neither locator reports: CodeMapper missed \wiFrIndCmQMiss\% of the behaviour-changing ones, the guarded version \wiFrIndCmGiMiss\%. AGi catches windows that land on an identical copy. The price is paid on preserved code: extra requests rose from \wiFrCmQFalse\% to \wiFrCmGiFalse\% on the stress seed, where most cases are adversarial, and from \wiStuCmQFalse\% to \wiStuCmGiFalse\% on the student repositories.

\textbf{The corrected copy guard costs less but does not dominate.} Over CodeMapper, AGi asked for fewer unnecessary re-checks than the original copy guard on every population (stress seed \wiFrCmGiFalse\% against \wiFrCmGFalse\%; student repositories \wiStuCmGiFalse\% against \wiStuCmGFalse\%; real histories \wiRealCmGiFalse\% against \wiRealCmGFalse\%). It missed more only where an identical copy sits just below the original, the case in which the line diff aligns the original's lines with the copy: on the duplicate suite \wiFrDupCmGiMiss\% against \wiFrDupCmGMiss\%, which accounts for the whole difference on the stress seed, and \wiHumCmGiMissK{} against \wiHumCmGMissK{} misses on the human-annotated regions. The registered hypothesis that AGi lowers unnecessary requests without losing safety therefore held on the held-out short-anchor and compound suites but not on the held-out seed as a whole. The provenance rule AG2 (Section~\ref{sec:guards}) removed more unnecessary requests but failed exactly where a copy guard is needed: on the held-out duplicate suite the cascade with AG2 missed \wiFrDupCasAGtwoMiss\% of the needed re-checks, against \wiFrDupCasAGiMiss\% with AGi and \wiFrDupDepMiss\% with no guard. We recommend AGi, keep the original guard where duplicates are expected, and report AG2 as a negative result.

\textbf{A language model as the locator.} Magic Markup re-anchors annotations with a language model \cite{misback2024magic}, so we also placed one inside the workflow. The protocol was registered before any evaluation call, and its amendments, including a change of model made before that model's first evaluation call, are in the artifact. The model is \wjModel{}, the model the deployed portal is configured with, at temperature~0. For each anchor it receives the old region with 15 lines of context and the complete new version of the paired file. It returns where the region now is and whether it changed. Its own answer can decide the request (\emph{own judgement}), or its window can pass through the same labels and guards as CodeMapper's (Table~\ref{tab:llm}). Only public code was sent: the human-annotated regions, the open-source edit-targeted histories and the student repositories.

% Generated by tools/v11/build_wpj_tex.py from summary-20261104T092825.json; do not edit.
\begin{table*}[t]
\caption{A Language-Model Locator With and Without the Safety Layer (\%)}\label{tab:llm}
\centering\footnotesize
\setlength{\tabcolsep}{3pt}
\begin{tabular}{@{}lrrrrr@{\hspace{8pt}}rrrrr@{\hspace{8pt}}rrrrr@{}}\toprule
 & \multicolumn{5}{c}{Human-annotated} & \multicolumn{5}{c}{Open-source, edit-targeted} & \multicolumn{5}{c}{Student, edit-targeted} \\
 & \multicolumn{5}{c}{$n{=}$200} & \multicolumn{5}{c}{$n{=}$2,454} & \multicolumn{5}{c}{$n{=}$1,857} \\
\cmidrule(lr){2-6}\cmidrule(lr){7-11}\cmidrule(lr){12-16}
Method & Place & Miss & Silent & Misdir & Extra & Place & Miss & Silent & Misdir & Extra & Place & Miss & Silent & Misdir & Extra \\ \midrule
\multicolumn{16}{@{}l}{\emph{Language model (\wjModel{})}} \\
\quad own judgment & 91.7 & 6.2 & 2.0 & 3.0 & 35.7 & 94.4 & 2.1 & 0.4 & 3.1 & 11.3 & 90.2 & 0.5 & 0.4 & 4.6 & 25.3 \\
\quad labels only & 91.7 & 6.2 & 1.5 & 3.5 & 37.1 & 94.4 & 0.4 & 0.3 & 3.2 & 10.5 & 90.2 & 0.3 & 0.3 & 4.7 & 25.9 \\
\quad +IG+AGi & 91.7 & 3.1 & 0.0 & 5.0 & 45.7 & 94.4 & 0.0 & 0.0 & 3.5 & 18.3 & 90.2 & 0.0 & 0.1 & 5.0 & 40.4 \\
\multicolumn{16}{@{}l}{\emph{Second model (Nemotron~3.5 Super)}} \\
\quad own judgment & 88.5 & 5.4 & 1.5 & 3.0 & 35.7 & 95.3 & 0.6 & 0.5 & 2.0 & 10.5 & 91.8 & 0.6 & 0.3 & 3.4 & 25.3 \\
\quad +IG+AGi & 88.5 & 2.3 & 0.0 & 4.5 & 44.3 & 95.3 & 0.0 & 0.0 & 2.4 & 18.9 & 91.8 & 0.0 & 0.1 & 3.7 & 40.4 \\
\multicolumn{16}{@{}l}{\emph{CodeMapper}} \\
\quad labels only & 93.6 & 3.1 & 0.0 & 7.5 & 35.7 & 92.2 & 0.4 & 0.3 & 5.5 & 10.1 & 83.2 & 0.3 & 0.3 & 10.6 & 29.2 \\
\quad +IG+AGi & 93.6 & 2.3 & 0.0 & 7.5 & 44.3 & 92.2 & 0.0 & 0.0 & 5.9 & 17.6 & 83.2 & 0.0 & 0.1 & 10.8 & 40.9 \\
\multicolumn{16}{@{}l}{\emph{Cascade}} \\
\quad deployed & 62.2 & 8.5 & 3.5 & 9.0 & 34.3 & 65.8 & 0.9 & 0.8 & 5.7 & 0.0 & 73.6 & 2.6 & 2.3 & 6.1 & 0.4 \\
\bottomrule
\end{tabular}
\tablenote{}{Same observations for every row. The language model (\wjModel{}, the model of the deployed portal; temperature 0) sees the old region with context and the complete new version of the paired file. \emph{Own judgment}: the model's answer of whether the region changed decides the request. \emph{Labels only} and \emph{+IG+AGi}: the model's window passes through the quotation comparison and the window guards, as for CodeMapper; placement is unchanged. Unanswerable prompts (file too large for the context, provider refusal) count as execution errors and request re-check. The second model received the same prompts and settings; its window is placed by its own answer, so its placement differs. Columns as in Table~\ref{tab:locator}.}
\end{table*}


\textbf{The model located well, but its own judgement of change was not a safe signal.} It placed \wjPubSelfPlace\% of the surviving open-source targets and \wjStuSelfPlace\% of the student ones. CodeMapper placed \wjPubCmQPlace\% and \wjStuCmQPlace\%, and the deployed cascade \wjPubDepPlace\% and \wjStuDepPlace\%. On the human-annotated regions CodeMapper placed more (\wjHumCmQPlace\% against \wjHumSelfPlace\%). Judging change by itself, the model still left needed re-checks unrequested: \wjPubSelfMissK{} of \wjPubSelfMissN{} on the open-source histories (\wjPubSelfMiss\%), \wjStuSelfMissK{} of \wjStuSelfMissN{} on the student repositories and \wjHumSelfMissK{} of \wjHumSelfMissN{} on the human-annotated regions. It also placed \wjPubSelfSilentK{}, \wjStuSelfSilentK{} and \wjHumSelfSilentK{} targets on wrong code without a request.

\textbf{The safety layer made it safe.} With IG+AGi the model left \wjPubGiMissK{} needed re-checks unrequested on the open-source histories, \wjStuGiMissK{} on the student repositories and \wjHumGiMissK{} on the human-annotated regions. Silent misplacements fell to \wjPubGiSilentK{}, \wjStuGiSilentK{} and \wjHumGiSilentK{}. Both registered hypotheses held on every population. The paired cluster intervals of the change in missed re-checks were \wjPubJOneLo{} to \wjPubJOneHi{}, \wjStuJOneLo{} to \wjStuJOneHi{} and \wjHumJOneLo{} to \wjHumJOneHi{} points (J1), and those for silent misplacements also excluded zero (J2). On the human-annotated regions the result rests on few discordant regions, and the exact paired tests were not significant ($p=\wjHumJOneP$ and $p=\wjHumJTwoP$). As a secondary check, Holm's adjustment over the \wjHolmTests{} exact J1 and J2 tests leaves only the open-source pair below 0.05 (adjusted $p$ of \wjPubJOneHolm{} and \wjPubJTwoHolm{}); the student tests, which also rest on few discordant targets, give \wjStuJOneHolm{} and \wjStuJTwoHolm{}. The cost fell on preserved code: extra requests rose from \wjPubSelfFalse\% to \wjPubGiFalse\% on the open-source histories and from \wjStuSelfFalse\% to \wjStuGiFalse\% on the student repositories. A second model run on all three populations, \wjNemoModel{}, judged change more cautiously than \wjModel{} on the open-source histories (\wjNemoPubSelfMissK{} against \wjPubSelfMissK{} missed re-checks of \wjPubSelfMissN{}) and placed about as many targets (\wjNemoPubSelfPlace\% of the open-source and \wjNemoStuSelfPlace\% of the student ones). Under IG+AGi it too left \wjNemoPubGiMissK{} needed re-checks unrequested on the open-source histories, \wjNemoStuGiMissK{} on the student repositories and \wjNemoHumGiMissK{} of \wjNemoHumSelfMissN{} on the human-annotated regions; J1 was \wjNemoPubJOneSupported{}, \wjNemoStuJOneSupported{} and \wjNemoHumJOneSupported{}, and J2 \wjNemoPubJTwoSupported{}, \wjNemoStuJTwoSupported{} and \wjNemoHumJTwoSupported{} on the three populations. The safety layer therefore did not depend on which model located the code, and a model with a better judgment of change narrowed, but did not close, the gap the layer closes. Two other models, run on the human-annotated regions only, moved in the same direction. Missed re-checks fell from \wjStepHumSelfMissK{} to \wjStepHumGiMissK{} of \wjStepHumSelfMissN{} for \wjStepModel{} (J1 \wjStepHumJOneSupported{}) and from \wjStepFiveHumSelfMissK{} to \wjStepFiveHumGiMissK{} for \wjStepFiveModel{} (J1 \wjStepFiveHumJOneSupported{}). Temperature~0 did not make the model deterministic: a repeat of the human-annotated run gave the same placement, range and label for \wjRepSameAll{} of \wjRepN{} regions. These are dated observations of free models served by third parties.

\textbf{Which locator to deploy.} CodeMapper placed more anchors than the deployed cascade on the held-out seed (\wiFrCmQPlace\% against \wiFrDepPlace\%), the human-annotated regions (\wiHumCmQPlace\% against \wiHumDepPlace\% of surviving regions; Section~\ref{sec:external} also counts correctly reported deletions, \cExtHumCmAcc\% against \cExtHumDepAcc\%) and the real histories (\wiRealCmQPlace\% against \wiRealDepPlace\%), but it loses anchors that move to another file. With the cross-file fallback it placed \wiRealChQPlace\% on the real histories and \wiFrChQPlace\% on the stress seed, the most of any method on both, and kept that placement under the guards with \wiStuChGiMissK{} missed re-checks of \wiStuChGiMissN{} on the student repositories; this combination was assembled after the outcomes had been read and is exploratory. The guarded cascade at the calibrated threshold of 0.55 came close on the real histories (\wiRealCalGiPlace\%) and placed the most on the student repositories (\wiStuCalGiPlace\%), so the ranking of locators depends on the population. The safety layer behaved alike under every locator, and that is the property a deployment needs: the locator can be replaced as better ones appear while the change labels, the two window guards and the issue states, and with them the safety argument, stay the same. The language model is such a candidate. On the same edit-targeted open-source and student rows it placed more targets than CodeMapper or the deployed cascade, and under the same layer it missed as few needed re-checks there as CodeMapper (\wjPubGiMissK{} and \wjStuGiMissK{}, against \wjPubCmGiMissK{} and \wjStuCmGiMissK{}). It costs a model call per anchor, with a median of \wjPubMedianS{}\,s through a free service.
